\documentclass{article}
\usepackage{amsmath,amssymb}
\title{Higher uniformity data for a native feedback evaluator}
\date{15 November 2026}
\begin{document}\maketitle
\paragraph{Status and purpose.} Constructive local presentation, not a new universal-property theorem. It isolates the data missing from the ordinary two-dimensional feedback axiom presentation. Global assembly, finite-tower joining and their higher compatibility must still be supplied and checked before claiming a global evaluator.

\paragraph{The actual witness categories.} For a fixed external pair let $\mathcal N$ be the already constructed native dialgebra bicategory. Its object $x$ is a state/body pair $(X,f)$. For $x=(X,f)$ and $y=(Y,g)$, its hom category $E(x,y)$ has objects $(h,\alpha)$ and morphisms the full admitted pairs $(\sigma,\Gamma)$, with the actual boundary frames and process witness. At normalized levels built from bicategorical $K_0$, this resolves to state two-cells and the required $K_0$ equation. At general higher raw inputs the actual $\Gamma$ cannot be replaced by that equation or forgotten. Identity, composition, associator and unitors are those of the independently proved native construction.

\paragraph{Local semantic input.} Let $\mathcal D$ be the target hom bicategory for the same external pair. Supply:
\begin{enumerate}
\item an object $b(x)$ of $\mathcal D$ for every state/body object $x$;
\item a functor $U_{x,y}:E(x,y)\to\mathcal D(b(x),b(y))$ for every pair, acting on the entire witness $(\sigma,\Gamma)$;
\item invertible two-cells of $\mathcal D$,
$c_{v,u}:U(v)U(u)\Rightarrow U(vu)$ and $i_x:1_{b(x)}\Rightarrow U(1_x)$, natural in the full admitted witnesses.
\end{enumerate}
Since $\mathcal D$ is a process hom bicategory, its objects, arrows and two-cells are respectively target process one-, two- and three-cells. The last line therefore requests target process three-cells, not equalities of the images of native arrows.

Require the pseudofunctor associativity and unit equations for these data. For composable $u,v,w$, with the displayed bracketing and conventional right-to-left composition, the associativity equation compares
\[
 c_{w,vu}(1*c_{v,u})\,a_{\mathcal D}
 \quad\hbox{and}\quad
 U(a_{\mathcal N})\,c_{wv,u}(c_{w,v}*1).
\]
Here $a_{\mathcal N}$ is the actual native associator, including its process witness. The unit equations likewise use the actual native unitors and the chosen $i_x$. No equation may be checked solely on the state components when the native hom retains a genuine process witness.

\paragraph{Constructive local conclusion.} These data define a pseudofunctor $A:\mathcal N\to\mathcal D$: its object function is $b$, its hom functors are $U$, and its composition/identity constraints are $c,i$. Functoriality of $U$ handles vertical witness composition; the displayed equations and their units handle horizontal composition. Conversely any such pseudofunctor determines precisely these data. This is the pseudofunctor presentation applied to an explicitly known witness category, not a claim that suitable semantic $U,c,i$ always exist.

Fixed-state closure is obtained by composing $A$ with the independently constructed inclusion $C_X$. Thus separate operations $B_X$ need not be postulated a second time. If a target doctrine already specifies $B_X$, it must provide a chosen pseudonatural identification with $A C_X$, including the actual identity-state frames; agreement on body objects alone is insufficient.

\paragraph{Remaining global work.} For a process target tricategory, repeat this local input for each external pair. A global evaluation also needs sequential assembly adjoint equivalences, their full native naturality, the trihomomorphism associativity/unit modifications and their equations. Directed parallel needs its own comparison and four-process compatibility. Iterated closure additionally needs joining comparisons natural on nested full witnesses and compatible with the already constructed ordered higher coherences. Those conditions are not obtained merely by the local pseudofunctor theorem. The accompanying Conditional\_Evaluation\_of\_Generated\_Native\_Towers.tex proves evaluation of the generated finite tower from independently specified target operations, not a pseudoalgebra for an unconstructed unrestricted higher state monad.
\end{document}
